\chapter{AI/BI Dashboards and Genie}
\index{AI/BI}\index{dashboards}\index{Genie}\index{natural-language analytics}%
\index{semantic model}\index{certified dataset}\index{self-service BI}%
\begin{objectives}
\begin{itemize}
    \item Build AI/BI dashboards with curated datasets, parameters, and cross-filters over governed Gold data.
    \item Configure a Genie space with instructions, example queries, and trusted assets that yield correct SQL.
    \item Curate a semantic layer: certified views, named measures, and documented dimensions.
    \item Publish and schedule dashboards with governed credentials.
    \item Diagnose and fix wrong Genie answers with instructions rather than data hacks.
    \item Design a self-service analytics layer where certified and exploratory analytics coexist safely.
\end{itemize}
\end{objectives}

\begin{crosscloud}
AI/BI is Databricks' \textbf{native BI layer}; Genie is its natural-language front end. Every
platform now pairs BI with an LLM assistant.

\vspace{2mm}
{\footnotesize
\begin{tabular}{@{}p{3.0cm}p{3.4cm}p{3.2cm}p{3.2cm}@{}}
\toprule
\textbf{Capability} & \textbf{Databricks} & \textbf{AWS} & \textbf{Azure / GCP / Fabric} \\
\midrule
BI surface & AI/BI dashboards & QuickSight & Power BI / Looker / Power BI (Fabric) \\
NL analytics & Genie spaces & Amazon Q in QuickSight & Copilot (Power BI) / Gemini (Looker) \\
Semantic steering & Genie instructions + certified tables & Q topics & Semantic model + Copilot prep / LookML \\
Compute behind BI & SQL warehouse & SPICE / direct query & Import / DirectQuery / Direct Lake \\
Scheduling/sharing & Dashboard subscriptions & Scheduled emails & Subscriptions / Looker schedules \\
Governance link & Unity Catalog lineage to tiles & Q/QuickSight permissions & Purview / Looker permissions \\
\bottomrule
\end{tabular}}

\vspace{1mm}
The differentiator to articulate in interviews: AI/BI reads the governed Delta tables in place and
inherits Unity Catalog lineage, so a dashboard tile traces back to the pipeline that produced it.
\end{crosscloud}

\section{Week 7 Roadmap and Mental Model}
Weeks 1--6 built governed, performant Gold data. Week 7 makes it \emph{consumable}: dashboards for the known questions and Genie for the unknown ones. The engineering content is not chart placement---it is the semantic discipline that makes self-service safe: certified datasets, named measures, and instructions that encode business definitions so that ``revenue'' means one thing everywhere \cite{databricksDashboards,databricksGenie}.

The mental model: \textbf{a dashboard is a cached contract; Genie is a compiler}. Dashboards execute your SQL on a schedule and render the results---trust comes from the SQL being right and the source being certified. Genie compiles natural language into SQL at question time---trust comes from the space being scoped to clean Gold tables with instructions that teach it your vocabulary. Both fail the same way: confidently, when pointed at ungoverned data.

\begin{definitionbox}
\textbf{An AI/BI dashboard} is a Databricks-native BI artifact: a set of datasets (SQL against governed tables) feeding visualizations with parameters and cross-filters, published on a schedule. \textbf{A Genie space} is a curated natural-language analytics surface: selected tables plus instructions, example queries, and trusted assets that steer generated SQL.
\end{definitionbox}

\begin{keyterms}
\textbf{Certified table/view}: a Gold asset marked as the governed source for a metric. \textbf{Instruction}: Genie steering text---definitions, join rules, filters. \textbf{Trusted asset}: a pre-verified query Genie can use as an authoritative answer. \textbf{Parameter}: a dashboard input bound across datasets. \textbf{Cross-filter}: a selection in one visualization filtering the others.
\end{keyterms}

\section{Monday: AI/BI Dashboards}
\subsection{Datasets and Visualizations}
A dashboard's \textbf{datasets} are SQL queries (or table references) executed on a SQL warehouse; \textbf{visualizations} render them; \textbf{parameters} thread user input through every dataset. Author datasets against Gold views---never Bronze---and prefer views that already encode the business grain so the dataset SQL stays thin and reviewable.

\begin{lstlisting}[language=SQL,caption={A parameterized dashboard dataset over a certified view.},label={lst:ch7-dataset}]
SELECT order_date, region, category,
       SUM(net_amount) AS revenue,
       COUNT(DISTINCT order_id) AS orders
FROM de_training.gold.v_sales_certified      -- certified, governed
WHERE order_date BETWEEN :start_date AND :end_date
  AND region IN (:regions)
GROUP BY order_date, region, category;
\end{lstlisting}

\subsection{Parameters, Cross-Filters, and Publishing}
Parameters make one dashboard serve many slices (dates, regions, channels). Publishing snapshots the dashboard on a schedule with the publisher's warehouse and credentials; subscribers receive rendered output. The governance consequence: \emph{whoever publishes shares their access}---publish with a service identity whose Gold read rights are exactly what subscribers should see (Chapter 8 makes this enforceable).

\section{Tuesday: Genie Spaces}
\subsection{What Genie Needs to Succeed}
Genie answers questions by generating SQL over the tables you give it, steered by natural-language context \cite{databricksGenie}. The inputs, in order of leverage: (1) few, clean, well-named Gold tables with documented columns; (2) \textbf{instructions} defining business terms (``revenue excludes cancelled orders''), join paths, and default filters; (3) \textbf{example queries} pairing questions with verified SQL; (4) \textbf{trusted assets} for the exact certified answers.

\begin{pitfall}
\textbf{Genie over Bronze.} Pointed at raw tables, Genie produces fluent, plausible, wrong SQL: unparsed dates, pre-dedup counts, business terms guessed from column names. Genie amplifies your data quality---it does not create it. Scope spaces to certified Gold only.
\end{pitfall}

\subsection{Debugging a Wrong Answer}
The Genie workflow when an answer is wrong: inspect the generated SQL, identify the semantic gap (wrong join? wrong filter? wrong measure?), then fix it \emph{in the space}---an instruction, an example query, or a trusted asset---never by editing data to match the answer. Each fix generalizes: a good instruction repairs a whole class of future questions.

\section{Wednesday: The Semantic Layer}
\subsection{Certified Views and Named Measures}
Between Gold tables and consumers sits a thin, deliberate semantic layer: views that rename columns to business language, encode the certified metric definitions, and hide operational columns. The goal is that every tool---dashboards, Genie, notebooks, external BI---computes ``revenue'' identically because they all read the same view.

\begin{lstlisting}[language=SQL,caption={A certified semantic view with named measures.},label={lst:ch7-certified}]
CREATE OR REPLACE VIEW de_training.gold.v_sales_certified AS
SELECT  o.order_date,
        c.region,
        p.category,
        o.net_amount,                       -- certified: net of cancels/refunds
        o.order_id
FROM    de_training.gold.fact_orders o
JOIN    de_training.gold.dim_customer c USING (customer_id)
JOIN    de_training.gold.dim_product  p USING (product_id)
WHERE   o.order_status NOT IN ('cancelled', 'refunded');
\end{lstlisting}

\begin{tipbox}
Certification is a process, not a flag: a view becomes certified when finance reconciles it to the source system, an owner is named, and changes go through review. Then mark it---and grant broadly.
\end{tipbox}

\section{Thursday Hands-on Project: Executive Dashboard and Genie Space}
Build a three-page executive dashboard on the certified Gold view, then a Genie space that correctly answers five business questions---including one it initially gets wrong.
\index{hands-on lab}\index{Genie!lab}\index{dashboard!lab}

\begin{lstlisting}[language=SQL,caption={Week 7 lab assets: certified view plus the dashboard's core datasets.},label={lst:ch7-lab}]
-- 0. Certified semantic view (from Wednesday)
CREATE OR REPLACE VIEW de_training.gold.v_sales_certified AS
SELECT o.order_date, c.region, p.category, o.net_amount, o.order_id
FROM de_training.gold.fact_orders o
JOIN de_training.gold.dim_customer c USING (customer_id)
JOIN de_training.gold.dim_product  p USING (product_id)
WHERE o.order_status NOT IN ('cancelled','refunded');

-- Dataset A: revenue trend      Dataset B: revenue by region
-- Dataset C: top categories     (all parameterized on :start_date/:end_date)

-- Genie space over de_training.gold (certified tables only) with instructions:
--  - "Revenue" = SUM(net_amount) from v_sales_certified, excluding cancels.
--  - "Last quarter" means the most recent completed calendar quarter.
--  - Join customer to orders via customer_id only.
-- Trusted asset: verified SQL for "revenue by state last quarter".
\end{lstlisting}

\subsection*{Deliverables}
\begin{itemize}
    \item A published three-page dashboard (trend, geography, category) with date and region parameters.
    \item A Genie space over Gold with instructions, two example queries, and one trusted asset.
    \item A short log: five test questions, the one wrong answer, its generated SQL, and the instruction that fixed it.
\end{itemize}

\subsection*{Acceptance Criteria}
\begin{itemize}
    \item Every dashboard dataset reads certified Gold views---no Bronze or ad-hoc tables.
    \item ``Revenue by state last quarter'' in Genie matches the dashboard's certified number exactly.
    \item The wrong-answer fix is an instruction or trusted asset, not a data change, and generalizes (a second phrasing of the same question now works).
    \item The published dashboard runs on a service identity with Gold-only read access.
\end{itemize}

\subsection*{Stretch Goals}
\begin{itemize}
    \item Add a cross-filter from the region map to the trend page.
    \item Measure Genie's first-try accuracy on 10 questions before and after instructions.
    \item Schedule a Monday 7 a.m.\ subscription of the executive page to a distribution list.
\end{itemize}

\section{Friday Use Case: Certified Numbers and Curious Users}
A retailer's finance team requires that any number leaving the analytics platform be certified and reconciled; meanwhile 300 business users want to ask Genie ad-hoc questions. Both must coexist without a shadow-Excel economy.
\index{self-service}\index{certification}\index{finance}

\subsection{Requirements and Assumptions}
Certified metrics change quarterly at most. Business questions are 80\% recurring themes, 20\% novel. Finance audits quarterly and currently distrusts anything not from their own extracts.

\begin{figure}[H]
    \centering
    \resizebox{\textwidth}{!}{%
    \begin{tikzpicture}[
        box/.style={rectangle, rounded corners, draw=primary, thick, fill=primary!6, align=center, minimum width=3.2cm, minimum height=1.0cm, font=\small\bfseries},
        cert/.style={rectangle, rounded corners, draw=codegreen!60!black, thick, fill=green!8, align=center, minimum width=3.0cm, minimum height=0.9cm, font=\small\bfseries},
        flow/.style={-{Latex[length=2mm]}, draw=primary, thick}
    ]
        \node[box] (gold) at (0,0) {Gold tables\\(facts, dims)};
        \node[cert] (views) at (4.6,0) {Certified views\\named measures};
        \node[box] (dash) at (9.2,1.3) {AI/BI dashboards\\(finance)};
        \node[box] (genie) at (9.2,-1.3) {Genie space\\(business users)};
        \node[box, draw=red!60!black, fill=red!5] (audit) at (13.6,0) {Lineage + audit\\quarterly review};
        \draw[flow] (gold) -- (views);
        \draw[flow] (views) -- (dash);
        \draw[flow] (views) -- (genie);
        \draw[flow, dashed] (dash) -- (audit);
        \draw[flow, dashed] (genie) -- (audit);
    \end{tikzpicture}%
    }
    \caption{One semantic layer feeding both certified dashboards and exploratory Genie, with lineage into the audit review.}
    \label{fig:ch7-selfservice}
\end{figure}

\subsection{Architecture Decisions}
\textbf{One semantic layer, two doors.} Both dashboards and Genie read the same certified views, so a number cannot diverge between the CFO's dashboard and a manager's Genie question. Genie's space is scoped to certified tables; its instructions encode the finance-approved definitions; trusted assets hold the board-level answers verbatim.

\textbf{Trust mechanics.} Every certified view carries an owner and a reconciliation record (finance signs quarterly). Dashboard tiles inherit Unity Catalog lineage, so an auditor can click from a number to the view to the pipeline. Ad-hoc exploration happens in a separate ``explore'' Genie space over non-certified Gold, clearly labeled---curiosity is allowed, export to finance decks is not.

\begin{pitfall}
\textbf{Certifying the dashboard instead of the data.} A ``certified dashboard'' over uncertified tables certifies nothing durable---the next pipeline change silently alters it. Certify the view, let any tool read it, and audit at the view layer.
\end{pitfall}

\subsection{Risks and Mitigations}
\begin{table}[H]
\centering
\small
\begin{tabular}{@{}p{4.6cm}p{7.6cm}@{}}
\toprule
\textbf{Risk} & \textbf{Mitigation} \\
\midrule
Genie confidently wrong on a metric & Golden question set re-run on every instruction change \\
Two ``revenue'' definitions circulate & One certified view; all tools read it; lineage proves it \\
Dashboard subscribers see too much & Publish under a least-privilege service identity \\
Explore space leaks into finance decks & Explicit labeling, separate space, export guardrails \\
Instruction debt accumulates unreviewed & Quarterly Genie space review with accuracy trend \\
\bottomrule
\end{tabular}
\caption{Week 7 scenario risks and mitigations.}
\label{tab:ch7-risks}
\end{table}

\section{Design Decision Guide}
\index{decision guide!Week 7}%
Self-service analytics succeeds or fails on semantic decisions, not chart choices.

\begin{table}[H]
\centering
\small
\begin{tabular}{@{}p{3.4cm}p{4.4cm}p{4.6cm}@{}}
\toprule
\textbf{Decision} & \textbf{Choose the first when} & \textbf{Choose the second when} \\
\midrule
Dashboard vs.\ Genie & Known recurring questions & Exploratory, phrased-in-the-moment questions \\
Certified view vs.\ ad-hoc dataset & Any number that leaves the room & Analyst's own scratch exploration \\
Instruction vs.\ trusted asset & A class of questions needs steering & One exact recurring question needs a fixed answer \\
Service identity vs.\ publisher identity & Anything with subscribers & Personal draft dashboards only \\
AI/BI vs.\ external BI tool & Governance colocation wins & Organizational BI standard mandates it \\
Refresh cadence: decision vs.\ maximum & Always decision cadence & Never the platform maximum by default \\
\bottomrule
\end{tabular}
\caption{Week 7 decision guide.}
\label{tab:ch7-decisions}
\end{table}

The unifying principle: one semantic layer, many doors. Every row above is a way of keeping the certified definition between the data and the decision-maker.

\section{Operational Guidance for Week 7}
\index{operational guidance!Week 7}%
\subsection{Performance and Reliability}
Track dashboard refresh duration and failure rate per dashboard; a tile that silently serves yesterday's data erodes trust faster than an honest error. For Genie, maintain the golden question set and re-run it after every instruction change---instruction edits are code changes with regression risk. Version the certified views (change control with reconciliation sign-off) because every tool downstream inherits their definitions.

\subsection{Security and Governance Checklist}
\begin{itemize}
    \item Genie spaces scoped to certified Gold only; explore spaces clearly labeled.
    \item Dashboard publication identity has exactly the access subscribers should inherit.
    \item Subscriptions reviewed quarterly---recipient lists drift as roles change.
    \item Certified-view changes require reconciliation evidence before merge.
\end{itemize}

\subsection{Troubleshooting Playbook}
\begin{table}[H]
\centering
\small
\begin{tabular}{@{}p{3.6cm}p{4.2cm}p{4.8cm}@{}}
\toprule
\textbf{Symptom} & \textbf{Likely cause} & \textbf{First action} \\
\midrule
Genie answers with wrong metric & Missing/ambiguous instruction & Inspect generated SQL; add instruction or trusted asset \\
Dashboard serves stale data & Refresh schedule or warehouse cold & Check refresh log and warehouse event history \\
Parameter not filtering & Dataset SQL missing the parameter binding & Rebind the parameter in the dataset, not the visual \\
Subscription not received & Publish identity lost access & Republish under the service identity; check delivery log \\
\bottomrule
\end{tabular}
\caption{Week 7 troubleshooting quick reference.}
\label{tab:ch7-trouble}
\end{table}

\section{Interview Q\&A}
\subsection{Concepts and Trade-offs}
\begin{enumerate}
    \item \textbf{Q: Why not just give analysts Power BI against Gold?}\\
    A: You can---the point of AI/BI is colocation with governance: lineage to the tile, no extract pipelines, and one semantic layer shared with Genie. External BI is a valid choice when the organization's standard demands it; the certified views remain the contract either way.
    \item \textbf{Q: How do you stop Genie from hallucinating metrics?}\\
    A: Scope to certified tables, encode definitions as instructions, provide example queries and trusted assets, and measure first-try accuracy on a fixed question set before rolling out.
    \item \textbf{Q: Dashboard performance is poor. Where do you look?}\\
    A: Dataset SQL (is it scanning the fact table instead of an MV or the certified view?), warehouse sizing/queueing (Chapter 6), and table layout (Chapter 5)---in that order.
    \item \textbf{Q: Who should a published dashboard run as?}\\
    A: A service identity whose data access equals what subscribers may see; publishing as a privileged user silently escalates every subscriber's effective access.
    \item \textbf{Q: How do you handle a metric definition change?}\\
    A: Change the certified view under review, re-reconcile with finance, bump the Genie instructions in the same change---one definition, one place, every tool updates together.
    \item \textbf{Q: What is the first thing you check when a certified number changes unexpectedly?}\\
    A: The view's lineage and the pipeline history---certified numbers change because upstream changed, and lineage turns ``who touched this'' into a query.
    \item \textbf{Q: Can Genie replace a semantic layer?}\\
    A: No---instructions steer generation, they do not enforce definitions. The certified views are the semantics; Genie is one of several doors onto them.
    \item \textbf{Q: How many Genie spaces should a domain have?}\\
    A: Few: one certified space per domain plus one explore space. Every additional space forks the instructions and example set---and the definitions they encode.
\end{enumerate}

\section{Further Reading}
The AI/BI dashboards \cite{databricksDashboards} and Genie \cite{databricksGenie} documentation are the product references. Revisit SQL warehouses for the compute underneath \cite{databricksSQLWarehouses} and Unity Catalog for the certification and lineage mechanics \cite{databricksUnityCatalog,databricksLineage}.

\section*{Key Takeaways}
\begin{itemize}
    \item Dashboards are cached contracts; Genie is a compiler---both inherit the quality of the data and semantics you give them.
    \item Certified views with named measures are the single source of metric truth for every tool.
    \item Fix wrong Genie answers with instructions and trusted assets---never by editing data.
    \item Published dashboards run as the publisher; publish with a least-privilege service identity.
    \item Self-service and certification coexist when both doors open onto the same semantic layer.
\end{itemize}

\section{Exercises}
\begin{enumerate}
    \item \textbf{(Conceptual)} Explain to finance why a Genie answer and a dashboard tile can be trusted to agree.
    \item \textbf{(Conceptual)} List three instruction types that steer Genie and give an example of each.
    \item \textbf{(Hands-on)} Complete the Thursday lab, including the wrong-answer fix log.
    \item \textbf{(Hands-on)} Add a second Genie space over uncertified Gold labeled ``explore only'' and document the guardrail difference.
    \item \textbf{(Hands-on)} Trace one dashboard tile's lineage to its source tables in the catalog explorer; capture it.
    \item \textbf{(Design)} Write the certification checklist (owner, reconciliation, review cadence) for a new certified view.
    \item \textbf{(Operations)} A certified view must change schema. Write the change plan covering dashboards, Genie, and downstream notebooks.
    \item \textbf{(Architecture)} When would you move a heavy dashboard dataset into a materialized view (Chapter 6) rather than tuning the SQL? Quantify.
\end{enumerate}
